% ============================================================
%  第四讲
% ============================================================
\section{第四讲}

\subsection{Hecke 函子}
\srcpage{11}

现在对一般约化群 \(G\)（在 \(k\) 上分裂且连通）定义 Hecke 性质。设 \(\Gd\) 是 \(\e\) 上的对偶约化群。

我们首先有 \(G\)-丛的概念。

\begin{example}\label{ex:g-bundles}
\(G\)-丛具体表现为：

\begin{itemize}
\item \(G=\mathrm{GL}_n\)：秩 \(n\) 向量丛 \(\cE\)；
\item \(G=\mathrm{SL}_n\)：秩 \(n\) 向量丛 \(\cE\)，并带有 \(\det\cE\simeq\cO\)；
\item \(G=\mathrm{PGL}_n\)：秩 \(n\) 向量丛 \(\cE\)，模去张量线丛；
\item \(G=\mathrm{O}_n\)：秩 \(n\) 向量丛 \(\cE\)，并带有对称非退化形式 \(\cE\otimes\cE\to\cO\)；
\item \(G=\mathrm{Sp}_n\)：秩 \(n\) 向量丛 \(\cE\)，并带有交错非退化形式 \(\cE\otimes\cE\to\cO\)；
\item \(\cdots\)
\end{itemize}
\end{example}

对任意概形 \(S\) 上的 \(G\)-丛 \(P_G\)，有右 \(t\)-正合的对称幺半函子
\yizhu{当 \(S\) 是导出概形时，下面的函子并不 \(t\)-正合；原文的「右 \(t\)-正合」只在
普通概形上成立。译文保留此限定。}{原文省略·A：原文脚注 4}
\[
\Rep(G)\to\QCoh(S),\qquad V\mapsto V_{P_G}.
\]
若 \(V\) 有限维，则 \(V_{P_G}\) 是有限秩向量丛。

我们还有 \(\Gd\)-局部系统的概念，它可以定义为右 \(t\)-正合的对称幺半函子
\yizhu{原文此处写 \(\Rep(G)\to\Shv(X)\)，但按上下文（\(\Gd\)-局部系统、\(V_\sigma\)）应为 \(\Rep(\Gd)\)。
这是原文的笔误，译文照录原文并加此注。}{原文笔误·A：与上一段 \(\Rep(G)\to\QCoh(S)\) 对照}
\[
\Rep(G)\to\Shv(X),\qquad V\mapsto V_\sigma.
\]
若 \(V\in\Rep(\Gd)\) 有限维，则 \(V_\sigma\) 是 lisse 的。

现在对任意 \(V\in\Rep(\Gd)\)，可以定义函子
\[
H_V\colon\Shv(\Bun_G)\to\Shv(\Bun_G\times X),
\]
它推广了上一讲定义的 Hecke 函子（即 \(\Gd=\mathrm{SL}_2\)、\(V\) 为标准表示的情形）。
定义所需的部件是下述等价。

\begin{theorem}[几何 Satake]\label{thm:satake}
对任意 \(x\in X\)，有典范幺半函子
\[
\Rep(\Gd)\to\Shv(L^+_xG\backslash L_xG/L^+_xG),\qquad V\mapsto S_V,
\]
其中 \(L_xG:=G(k((t_x)))\)，\(L^+_xG:=G(k[[t_x]])\)，\(t_x\) 是 \(x\) 附近的局部坐标。
\end{theorem}

\begin{remark}\label{rem:satake-indep}
可以验证 \(S_V\) 不依赖于局部坐标 \(t_x\) 的选取。
\end{remark}

双商
\[
\mathrm{Hecke}^{\mathrm{loc}}_x:=L^+_xG\backslash L_xG/L^+_xG
\]
分类 \(D_x:=\operatorname{Spec} k[[t_x]]\) 上的两个 \(G\)-丛 \(P_G,P'_G\)，以及它们在
\(\overset{\circ}{D}_x\) 上的限制之间的一个同构。我们还可以考虑这个叠的整体版本
\(\mathrm{Hecke}_x\)，它分类 \(X\) 上的两个 \(G\)-丛，以及它们在
\(\overset{\circ}{X}:=X\setminus x\) 上的限制之间的一个同构。对应为
\[
\begin{tikzcd}
& \mathrm{Hecke}_x \arrow[dl,"\overleftarrow{h}"'] \arrow[dr,"\overrightarrow{h}"] \arrow[d,"\pi"] & \\
\Bun_G & \mathrm{Hecke}^{\mathrm{loc}}_x & \Bun_G
\end{tikzcd}
\]

\begin{definition}\label{def:hecke-functor}
定义函子
\[
H_{V,x}\colon\Shv(\Bun_G)\to\Shv(\Bun_G),\qquad
F\mapsto\overrightarrow{h}_*\bigl(\overleftarrow{h}^!(F)\overset{!}{\otimes}\pi^!(S_V)\bigr).
\]
让 \(x\in X\) 变动，类似地定义函子
\[
H_V\colon\Shv(\Bun_G)\to\Shv(\Bun_G\times X),
\]
使得 \(H_V(F)|_{\Bun_G\times\{x\}}\simeq H_{V,x}(F)\)。我们称它们为 \emph{Hecke 函子}。
\end{definition}

\begin{exercise}\label{ex:hecke-pgl2}
取 \(G=\mathrm{PGL}_2\)、\(\Gd=\mathrm{SL}_2\)，并考虑标准表示 \(\mathrm{std}\in\Rep(\Gd)\)。
验证上述 Hecke 函子 \(H_{\mathrm{std}}\) 等价于 \eqref{eq:3.2} 中的函子。
\end{exercise}

Hecke 函子以下述意义交换。对 \(V_1,V_2\in\Rep(\Gd)\)，令 \(H_{V_1,V_2}\) 为复合
\[
\Shv(\Bun_G)\xrightarrow{\ H_{V_2}\ }\Shv(\Bun_G\times X)
\xrightarrow{\ H_{V_1}\ }\Shv(\Bun_G\times X\times X).
\]
则有典范识别
\[
H_{V_1,V_2}\simeq\mathrm{swap}\circ H_{V_2,V_1},
\]
其中 \(\mathrm{swap}\) 是由交换两个 \(X\) 因子诱导的 \(\Shv(\Bun_G\times X\times X)\) 上的对合。
更一般地，给定有限集 \(I\) 与以 \(I\) 为指标集的 \(\Gd\)-表示族 \(V=\{V_i\}_{i\in I}\)，
有良定义函子
\[
H_V\colon\Shv(\Bun_G)\to\Shv(\Bun_G\times X^I),
\]
它由任意顺序复合 \(H_{V_i}\) 得到。

\begin{exercise}\label{ex:hecke-compose}
验证
\[
H_{V_1,x}\circ H_{V_2,x}\simeq H_{V_1\otimes V_2,x}.
\]
\end{exercise}

\begin{question}\label{ques:action}
什么作用在 \(\Shv(\Bun_G)\) 上？
\end{question}

对每个 \(x\in X\)，有 \(\Shv(\Bun_G)\) 上的自函子 \(H_{V,x}\)，但 \(H_V\) 并不是自函子。
有一种形式化的办法，可以从这些 Hecke 函子产生自函子。对有限集 \(I\)、
\(V=\boxtimes_{i\in I}V_i\in\Rep(\Gd)^{\otimes I}\) 以及 \(G\in\Shv(X^I)\)，考虑复合
\begin{equation}\label{eq:4.1}
\Shv(\Bun_G)\xrightarrow{\ H_V\ }\Shv(\Bun_G\times X^I)
\xrightarrow{\ -\overset{!}{\otimes}\mathrm{pr}_2^!(G)\ }\Shv(\Bun_G\times X^I)
\xrightarrow{\ \mathrm{pr}_{1,*}\ }\Shv(\Bun_G).
\end{equation}
这给出函子
\[
\Rep(\Gd)^{\otimes I}\otimes\Shv(X^I)\to\End(\Shv(\Bun_G)),
\]
它把 \(V\boxtimes G\) 送到上述自函子。

\begin{remark}\label{rem:integration}
当 \(I=\{1,2\}\) 时，自函子 \eqref{eq:4.1} 是函子 \(H_{V_1,x_1}\circ H_{V_2,x_2}\)（当 \(x_1\neq x_2\)）
与 \(H_{V_1\otimes V_2,x}\)（当 \(x_1=x_2=x\)）对 \(X^2\) 上的「测度」\(G\) 的「积分」。
\end{remark}

对有限集之间的任意映射 \(\alpha\colon I\to J\)，有函子
\[
\Rep(\Gd)^{\otimes I}\xrightarrow{\ \mathrm{mult}_\alpha\ }\Rep(\Gd)^{\otimes J},
\]
它由 \(\Rep(\Gd)\) 的对称幺半结构给出；还有函子
\[
\Delta^\alpha_*\colon\Shv(X^J)\to\Shv(X^I),
\]
它由映射 \(\Delta_\alpha\colon X^J\to X^I\) 诱导。可以验证下图交换：
\[
\begin{tikzcd}
\Rep(\Gd)^{\otimes I}\otimes\Shv(X^J) \arrow[r,"\mathrm{mult}_\alpha\otimes\mathrm{id}"] \arrow[d,"\mathrm{id}\otimes\Delta^\alpha_*"'] & \Rep(\Gd)^{\otimes J}\otimes\Shv(X^J) \arrow[d] \\
\Rep(\Gd)^{\otimes I}\otimes\Shv(X^I) \arrow[r] & \End(\Shv(\Bun_G)).
\end{tikzcd}
\]
于是有函子
\begin{equation}\label{eq:4.2}
\Rep(\Gd)_{\mathrm{Ran}}:=\operatorname*{colim}_{\alpha\colon I\to J}
\Rep(\Gd)^{\otimes I}\otimes\Shv(X^J)\to\End(\Shv(\Bun_G)),
\end{equation}
其中余极限按映射 \(\alpha\colon I\to J\) 取，并对 \(I\) 反变、对 \(J\) 协变。

注意 \(\Rep(\Gd)_{\mathrm{Ran}}\) 带有（对称）幺半结构：把 \(I\xrightarrow{\alpha}J\) 与
\(I'\xrightarrow{\alpha'}J'\) 的项张量成 \(I\sqcup I'\xrightarrow{(\alpha,\alpha')}J\sqcup J'\) 的项。
可以验证 \eqref{eq:4.2} 与两边的幺半结构相容。
\yizhu{\(\Rep(\Gd)_{\mathrm{Ran}}\) 的「Ran」指 Ran 空间：它把 \(X\) 的所有有限子集放在一起，
允许这些点碰撞。\(\Shv(X^I)\) 的因子记录 Hecke 修改发生的位置；当两个点碰撞时，两个 Hecke 函子复合，
对应 \(\Rep(\Gd)\) 的张量积。这个构造的用途在 \S4.2 显现：它让所有 Fourier 系数有同一个「家」。}{背景直觉·A：Gaitsgory--Raskin, GLC I, arXiv:2405.03599, \S2（\(\Rep(\Gd)_{\mathrm{Ran}}\) 在 \(\Shv(\Bun_G)\) 上的作用）}

总之，把所有 Hecke 函子放在一起，得到幺半函子
\[
\Rep(\Gd)_{\mathrm{Ran}}\to\End(\Shv(\Bun_G)).
\]

\subsection{Hecke 性质}
\srcpage{14}

\begin{question}\label{ques:ran}
\(\Rep(\Gd)_{\mathrm{Ran}}\) 能为我们做什么？
\end{question}

我们给出两个答案：

\begin{itemize}
\item 它把所有 Fourier 系数收进同一个「家」；
\item 它给出特征层的完整而正确的定义。
\end{itemize}

先解释第二点。固定 \(X\) 上的 \(\Gd\)-局部系统 \(\sigma\)，有对称幺半函子
\[
\mathrm{ev}_\sigma\colon\Rep(\Gd)_{\mathrm{Ran}}\to\Vect,
\]
它在 \(I\xrightarrow{=}I\) 这一项上的限制为
\[
\Rep(\Gd)^{\otimes I}\otimes\Shv(X^I)
\xrightarrow{\ \sigma^{\boxtimes I}\otimes\mathrm{id}\ }
\Shv(X^I)\otimes\Shv(X^I)
\xrightarrow{\ -\overset{!}{\otimes}-\ }
\Shv(X^I)\xrightarrow{\ \Gamma\ }\Vect,
\]
其中回忆：\(X^I\) 上的 \(\Gd\)-局部系统 \(\sigma^{\boxtimes I}\) 被看作对称幺半函子
\(\Rep(\Gd)^{\otimes I}\to\Shv(X^I)\)。

\begin{definition}\label{def:eigensheaf}
以 \(\sigma\) 为特征值的 \emph{Hecke 特征层}是一个 \(\Rep(\Gd)_{\mathrm{Ran}}\)-线性函子
\[
\Vect\to\Shv(\Bun_G),\qquad \one\mapsto F,
\]
其中 \(\Rep(\Gd)_{\mathrm{Ran}}\) 经由 \(\mathrm{ev}_\sigma\) 作用在 \(\Vect\) 上。
我们常常滥用记号，把 Hecke 特征层本身就记作 \(F\)。
\end{definition}

\begin{exercise}\label{ex:sigma-unique}
证明：\(\sigma\) 由对称幺半函子 \(\mathrm{ev}_\sigma\) 唯一确定。
\end{exercise}

\begin{exercise}\label{ex:eigen-hecke}
设 \(F\) 是以 \(\sigma\) 为特征值的 Hecke 特征层。证明
\begin{equation}\label{eq:4.3}
H_V(F)\simeq F\boxtimes V_\sigma.
\end{equation}
\end{exercise}

\begin{remark}\label{rem:generalize}
注意 \eqref{eq:4.3} 推广了 \eqref{eq:3.3}。
\end{remark}

\subsection{几何 Langlands 对应}
\srcpage{14}

\begin{conjecture}\label{conj:gl}
若 \(\sigma\) 是不可约 \(\Gd\)-局部系统（即它不来自真抛物子群），则存在唯一的 Hecke 特征层 \(\cF_\sigma\)，
并带有同构 \(\coeff_0(\cF_\sigma)\simeq\one\)。
\end{conjecture}

\begin{theorem}[Arinkin 等]\label{thm:gl-general}
{\footnotesize\noindent 完整作者名单：Arinkin--Beraldo--Campbell--Chen--Faergeman--Gaitsgory--Lin--Raskin--Rozenblyum（原文以 ``R.'' 代 Raskin）。\par}
若 \(\operatorname{char}(k)=0\)，则猜想 \ref{conj:gl} 成立。此外还有：
\begin{itemize}
\item \(\cF_\sigma\) 是反常的；
\item \(\cF_\sigma\) 是尖点的；
\item \(\cF_\sigma\) 是半单的。更精确地说，设 \(S_\sigma\) 是 \(\sigma\) 的自同构群，则有分解
\[
\cF_\sigma\simeq\bigoplus_{\rho\in\mathrm{Irr}(S_\sigma)}\cF_{\sigma,\rho}^{\oplus\dim\rho},
\]
其中 \(\cF_{\sigma,\rho}\) 是单反常层；
\item \(\cF_\sigma\) 的特征圈是 \([\mathrm{Nilp}]\)，其中 \(\mathrm{Nilp}\subset T^*\Bun_G\) 是整体幂零锥；
\yizhu{特征圈为 \([\mathrm{Nilp}]\) 说明 \(\cF_\sigma\) 的奇异支撑是整个幂零锥 \(\mathrm{Nilp}\subset T^*\Bun_G\)。
这是整体几何 Langlands 的深刻结论，与尖点性并不矛盾：尖点性是关于常数项消没的，而奇异支撑描述的是
微局部行为。}{背景直觉·A：Gaitsgory--Raskin, GLC I, arXiv:2405.03599, §6}
\item 当 \(g>1\) 时，\(\cF_\sigma\) 的通用秩为
\[
\prod_i d_i^{(2d_i-1)(g-1)},
\]
其中 \(d_i\) 是 \(G\) 的指数。
\yizhu{原文称 \(d_i\) 为 \(G\) 的指数。按此约定，\(\mathrm{GL}_2\) 的指数为 \(1,2\)，公式给出
\(1^{1}\cdot 2^{3(g-1)}=2^{3g-3}\)，与已知的 \(\mathrm{GL}_2\) Hecke 特征层通用秩一致；
\(\mathrm{PGL}_2\) 的指数为 \(1\)，公式给出 \(1\)，也与其特征层通用秩一致。因此原文公式无需改动。
注意不要把 \(d_i\) 读成不变多项式的次数：\(\mathrm{PGL}_2\) 的不变多项式次数是 \(2\)，代入会得到
不同的数值，那不是原文的约定。}{术语背景·A：原文公式；Gaitsgory--Raskin, GLC I, arXiv:2405.03599, §6}
\end{itemize}
\end{theorem}

\begin{remark}\label{rem:components}
\(\Bun_G\) 的连通分支可以与 \(\mathrm{Irr}(Z(\Gd))\)（\(\Gd\) 的中心的不可约表示之集合）等同。
反常层 \(\cF_{\sigma,\rho}\) 支撑在由 \(\rho|_{Z(\Gd)}\) 标号的连通分支上。
\end{remark}
